\scanpage{22}
\mspage{20}

\margin{dans $\mathrm{Hot}_{\bullet}$.}

pour faire le lien avec la théorie \unsure{primitive|première|ponctuée}, ou théorie générale \unsure{des|dans} $\mathrm{Hot}_{\bullet}$. Si on désigne par

\[
 \gamma:\mathrm D_{\geq 0}(k)\mathbin{\unsure{-|\longrightarrow}}\mathrm{Hot}_{\bullet}
\]

le premier de ces foncteurs, il ne commute pas aux produits $*$ ($\overset{\unsure{\mathbb L|L}}{\otimes}_k$ et $\times$ \unsure{resp.$^t$}), mais on a, pour $P_\bullet,Q_\bullet$ variables dans $\mathrm D_{\geq}(k)$, une flèche fonctorielle canonique

\[
 \gamma(P_\bullet\overset{\mathbb L}{\otimes}_k Q_\bullet)
 \xleftarrow{\mathrm{can.}}\underset{\text{``$k$-bilinéaire''}}{\gamma(P)\times\gamma(Q)}
\]

Ce qui montre que de donner des accouplements

\[
 P_{\unsure{\bullet}}\mathbin{*}Q\xrightarrow{\unsure{c_k|c_K}}R,
\]

d'où

\[
 \gamma(P\mathbin{\unsure{*|\circ}}Q)\xrightarrow{\gamma(\unsure{c_k|\varphi})}\gamma(R),\tag*{$(*)$}
\]

est l'un et l'autre plus fin que de se donner seulement

\[
 \gamma(P)\mathbin{*}\gamma(Q)\xrightarrow{c_\bullet}\gamma(R)\tag*{$(*')$}
\]

vu que \inserted{une donnée} $(*)$ implique une donnée $(*')$ par

\diagram{\begin{tikzcd}[ampersand replacement=\&,row sep=small]
 \&\gamma(P*Q)\arrow[dr,"{\gamma(c)}"]\&\\
 \gamma(P)\times\gamma(Q)\arrow[ur,"\mathrm{can.}"]\arrow[rr,"{c_\bullet}"']\&\&\gamma(R).
\end{tikzcd}}

Ceci vu, on s'attend bien à avoir une théorie \uline{forte} des accouplements
% archive sheet number: 21
